\documentclass[10pt,letterpaper]{article}
\usepackage[spanish]{babel}
\usepackage{listings}
\usepackage{graphicx} %paqute para que pueda jalar mi imagen del diagrama
\usepackage{float}
\usepackage{comment}%para comentar desde la pergunta 11 a la 15
% Preámbulo

\title{Documentacion de un Sistema de Gestion para una Tienda Online}
\author{Beatriz Garcia Lopez }
\date{\today}

\usepackage{xcolor}

\definecolor{codegreen}{rgb}{0,0.6,0}
\definecolor{codegray}{rgb}{0.5,0.5,0.5}
\definecolor{codepurple}{rgb}{0.58,0,0.82}
\definecolor{backcolour}{rgb}{0.95,0.95,0.92}

\lstdefinestyle{mystyle}{
    backgroundcolor=\color{backcolour},   
    commentstyle=\color{codegreen},
    keywordstyle=\color{magenta},
    numberstyle=\tiny\color{codegray},
    stringstyle=\color{codepurple},
    basicstyle=\ttfamily\footnotesize,
    breakatwhitespace=false,         
    breaklines=true,                 
    captionpos=b,                    
    keepspaces=true,                 
    numbers=left,                    
    numbersep=5pt,                  
    showspaces=false,                
    showstringspaces=false,
    showtabs=false,                  
    tabsize=2
}
\lstset{style=mystyle}

\begin{document}
\maketitle




\section{Introducción}
Se nos ha 
\section{Descripción del problema}
Lo que hace 
\section{Análisis y Diseño}
En este apartado vamos a estructurar un poco las tablas y implementar un DER(Diagrama Entidad-Relacion) utilizando la página "Draw.io"
\begin{table}[H]%para que lo ponga justo abajo de "Analisis"
    \begin{tabular}{| c |p{9cm} |} % esto sirve para dar un salto de linea para que no se siga de largo en LaTeX
        \hline
        TABLAS & ¿QUE ORGANIZA? \\
        \hline
        Categorias                 & Organiza los productos en 5 categorias diferentes
                                      (telefonos, laptops, accesorios, tablets y relojes inteligentes)\\ \hline
        
        Reseñas                    & Registra las reseñas de los productos y el ranking del mismo, el cual se evalua de
                                      1 a 5 estrellas, siendo 1 = muy malo y 5 = excelente. Un detalle importante es que
                                      las reseñas solo podran ser realizadas por clientes que hayan comprado un producto,
                                      sin importar la categoria del mismo.\\ \hline

        Pedidos                    & Especifica los detalles especificos del producto; como fecha de pedido y el estado
                                      del mismo (Enviado, Pendiente y Entregado).\\ \hline
        
        Detalles del Pedido        & Guarda los datos especificos del pedido como por ejemplo
                                      la cantidad (¿cuantos productos se van a enviar?) y el precio unitario del producto.\\ \hline
        
        Productos                  & Esta tabla es la mas importante ya que desde aqui vamos a poder gestionar todas las demas tablas.
                                      Es como un "gestor que se conecta con todo lo demas". De igual manera, aqui tenemos el nombre, descripcion, precio y stock de todos los
                                      productos disponibles al momento.\\ \hline

        Clientes                   & En esta tabla esta toda la informacion personal del cliente como por ejemplo nombre, correo, telefono y direccion.\\ \hline

    \end{tabular}
\end{table}

\section{Implementacion de la Base de Datos}

\section{Consultas y Procedimientos Almacenados}

\section{Validacion y Optimización}



\end{document}